% Generated draft: requires mathematical and source review.
\documentclass[11pt]{article}
\usepackage[margin=25mm]{geometry}
\usepackage{fontspec}
\setmainfont{Latin Modern Roman}
\usepackage{amsmath,amssymb,mathtools,unicode-math}
\setmathfont{Latin Modern Math}
\usepackage{enumitem,needspace,xurl,hyperref}
\setlength{\parindent}{0pt}
\setlength{\parskip}{5pt}
\setlength{\emergencystretch}{3em}
\newcommand{\N}{\mathbb{N}}
\newcommand{\Z}{\mathbb{Z}}
\newcommand{\Q}{\mathbb{Q}}
\newcommand{\R}{\mathbb{R}}
\newcommand{\C}{\mathbb{C}}
\newcommand{\F}{\mathbb{F}}
\newcommand{\A}{\mathbb{A}}
\newcommand{\PP}{\mathbb{P}}
\newcommand{\E}{\mathbb{E}}
\newcommand{\eps}{\varepsilon}
\newcommand{\dd}{\,\mathrm{d}}
\begin{document}
\section*{UnsolvedMath ID: 30000541}
\textbf{Source problem code:} catalogue entry
\subsection{Definitions and mathematical statement}
For each integer $n\ge3$, let $\mu$ be the intended smooth volume form on $S^n$. Consider the group of volume-preserving diffeomorphisms isotopic to the identity, conventionally written
\[
 G_n=\operatorname{Diff}_0(S^n,\mu).
\]
Volume preservation means $f^*\mu=\mu$. The source does not say whether the isotopy must stay volume preserving or only be an ambient diffeomorphism isotopy, nor does it state the differentiability/topology convention. The standard smooth identity-component interpretation must be checked.

Is $G_n$ bounded or unbounded for $n\ge3$? The meaning of boundedness is not defined in the supplied statement. In a common group-theoretic interpretation, a conjugation-invariant norm is a finite-valued function $\nu:G_n\to[0,\infty)$ satisfying
\[
 \nu(g)=0\Longleftrightarrow g=1,\quad
 \nu(g^{-1})=\nu(g),\quad
 \nu(gh)\le\nu(g)+\nu(h),\quad
 \nu(hgh^{-1})=\nu(g).
\]
Under that interpretation boundedness would mean
\[
 \forall\nu\text{ of this kind}:\quad\sup_{g\in G_n}\nu(g)<\infty,
\]
and unboundedness would mean existence of one such $\nu$ with infinite supremum. The bound may depend on $\nu$ and $n$; no single constant uniform over all norms or dimensions is required.

This norm-based interpretation is explanatory, not a supplied definition. If the report instead fixes a particular metric, restricts to continuous norms, or uses another boundedness property, the question changes. The source gives no such choice. The requested task should determine the behavior for each $n\ge3$, allowing any dimension dependence, rather than merely exhibit one dimension with a stated behavior.
\subsection{Short English statement}
The question asks whether the identity-isotopic volume-preserving diffeomorphism group of each sphere of dimension at least three is bounded. A usual interpretation uses all conjugation-invariant norms, but the source must specify boundedness and the isotopy convention.
\subsection{Sources}
\begin{itemize}
\item Local UnsolvedMath registry, permanent ID 30000541, code catalogue entry, ``Boundedness of Volume-Preserving Diffeomorphism Groups of Spheres.'' Proposed year 2006; proposer unspecified.
\item \emph{Geometric Group Theory, Hyperbolic Dynamics and Symplectic Geometry} (2007), Oberwolfach Reports, pp.\ 1991--2058, DOI 10.4171/owr/2006/33; \url{https://doi.org/10.4171/owr/2006/33}. The displayed report year 2007 differs from the DOI and proposed year 2006.
\item Statement status corrected\_verified; the supplied note says editorial and LaTeX cleanup preserved meaning and added no source-dependent content.
\item Literature triage checked 2026-08-21: open. It describes the boundedness question for volume-preserving diffeomorphisms of $S^n$, $n\ge3$, as unresolved in this generality, without specifying a norm or metric. Created 2026-08-25; updated 2026-09-16. The intended property and current status require independent verification.
\item Registry curation by Ulam AI / UnsolvedMath Contributors, CC BY 4.0; the underlying report retains its supplied provenance. Unreviewed statement-only draft.
\end{itemize}
\end{document}
